\documentclass{article}
\usepackage[utf8]{inputenc}
\usepackage{amsmath}
\usepackage{amssymb}
\usepackage{geometry}
\geometry{a4paper, margin=1in}
\title{Performance Loss in TDOA Estimation with 1-Bit Quantization}
\author{}
\date{}

\begin{document}
	
	\maketitle
	
	\section*{Performance Loss in TDOA Cross-Correlation Estimation with 1-Bit Quantization}
	
	\section*{Method 1: Based on Taylor Series Expansion}
	
	\subsubsection*{Objective}
	Due to the non-linear mapping relationship described by the Arcsine Law, TDOA estimates $\tau_q$ based on the quantized measurement cross-correlation $R_q(\tau)$ exhibit bias. Let $\hat{\tau}$ be the TDOA estimate obtained from the original cross-correlation $\rho(\tau)$, and $\tau_0$ be the true TDOA value. Then $\hat{\tau}$ is an unbiased estimate of $\tau_0$. Define $\Delta \tau_q = \tau_q - \tau_0$ and $\Delta \tau = \hat{\tau} - \tau_0$ as the corresponding estimation errors.
	
	Assumption: $\tau_q$ is a biased estimate of $\tau_0$, i.e.,
	\begin{gather}
		\tau_q = \hat{\tau} + \delta_\tau  \\
		\mathbb{E}[\tau_q] = \mathbb{E}[\hat{\tau}] + \mathbb{E}[\delta_\tau] = \tau_0 + \mathbb{E}[\delta_\tau] 
	\end{gather}
	where $\delta_\tau$ is the bias term and $\mathbb{E}[\delta_\tau] \propto \text{Var}(\hat{\tau})$. The research objective is to derive an expression for $\delta_\tau$.
	
	\subsubsection*{A. Model Assumptions}
	The signal is a known waveform $s(t)$. The received signals are
	\begin{equation}
		\begin{aligned}
			x(t) &= s(t) + n_1(t) \\
			y(t) &= s(t - \tau) + n_2(t)
		\end{aligned}
	\end{equation}
	
	where $n_1(t)$ and $n_2(t)$ are zero-mean Gaussian white noise processes. The goal is to estimate the delay $\tau$ over the observation interval $t \in [0, T]$.
	
	\subsubsection*{B. Cross-Correlation under 1-Bit Quantization: Arcsine Law}
	Apply 1-bit quantization to $x(t)$ and $y(t)$
	\begin{align}
		x_q(t) &= \text{sign}(x(t)), \quad y_q(t) = \text{sign}(y(t)) 
	\end{align}
	The cross-correlation of the quantized signals is
	\begin{align}
		R_q(\tau) &= \mathbb{E}[x_q(t)y_q(t - \tau)] 
	\end{align}
	If $x(t)$ and $y(t)$ are continuous-time, scalar, stationary real Gaussian processes, then according to the Arcsine Law (Van Vleck Formula)
	\begin{align}
		R_q(\tau) &= \frac{2}{\pi} \arcsin(\rho(\tau)) 
	\end{align}
	where $\rho(\tau)$ is the normalized correlation coefficient of the original signals
	\begin{align}
		\rho(\tau) &= \frac{R_{xy}(\tau)}{\sigma_x\sigma_y} = \frac{R_{xy}(\tau)}{\sqrt{R_{xx}(0)R_{yy}(0)}} 
	\end{align}
	Then the first and second order derivatives of $R_q(\tau)$ are
	
	\begin{align}
		R_{q}^{\prime}(\tau)&=\frac{2}{\pi} \cdot \frac{\rho^{\prime}(\tau)}{\sqrt{1-\rho^{2}(\tau)}} \\
		R_{q}^{\prime \prime}(\tau)&=\frac{2}{\pi} \cdot\left(\frac{\rho^{\prime \prime}(\tau)}{\sqrt{1-\rho^{2}(\tau)}}+\frac{\rho(\tau)\left[\rho^{\prime}(\tau)\right]^{2}}{\left(1-\rho^{2}(\tau)\right)^{3 / 2}}\right)
	\end{align}
	
	\subsubsection*{C. Taylor Expansion}
	Let $\tau_q$ be the TDOA estimate obtained from $R_q(\tau)$, $\hat{\tau}$ be the estimate obtained from $\rho(\tau)$, and $\tau_0$ be the true TDOA. Define $\Delta\tau_q = \tau_q - \tau_0$ and $\Delta\tau = \hat{\tau} - \tau_0$ as the estimation errors.
	
	\begin{itemize}
		\item Expand $R_q(\tau)$ in a Taylor series around $\hat{\tau}$
		\begin{equation}
			\begin{aligned}
				R_q(\tau) &= R_q(\hat{\tau}) + R_q'(\hat{\tau})(\tau - \hat{\tau}) + \frac{1}{2} R_q''(\hat{\tau})(\tau - \hat{\tau})^2  \\
				&= \frac{2}{\pi} \arcsin(\rho(\hat{\tau})) + \frac{2}{\pi} \cdot \frac{\rho'(\hat{\tau})}{\sqrt{1 - \rho^2(\hat{\tau})}} \cdot (\tau - \hat{\tau})  \\
				&\quad + \frac{1}{2} \cdot \frac{2}{\pi} \left( \frac{\rho''(\hat{\tau})}{\sqrt{1 - \rho^2(\hat{\tau})}} + \frac{\rho(\hat{\tau})[\rho'(\hat{\tau})]^2}{(1 - \rho^2(\hat{\tau}))^{3/2}} \right) \cdot (\tau - \hat{\tau})^2  \\
				&= \frac{2}{\pi} \arcsin(\rho(\hat{\tau})) + \frac{1}{2} \cdot \frac{2}{\pi} \cdot \frac{\rho''(\hat{\tau})}{\sqrt{1 - \rho^2(\hat{\tau})}} \cdot (\tau - \hat{\tau})^2 
			\end{aligned}
		\end{equation}
		
		Substituting $\tau_q$ gives
		\begin{equation}
			\begin{aligned}
				R_q(\tau_q) &= R_q(\hat{\tau}) + R_q'(\hat{\tau})(\tau_q - \hat{\tau}) + \frac{1}{2} R_q''(\hat{\tau})(\tau_q - \hat{\tau})^2  \\
				&= \frac{2}{\pi} \arcsin(\rho(\hat{\tau})) + \frac{1}{2} \cdot \frac{2}{\pi} \cdot \frac{\rho''(\hat{\tau})}{\sqrt{1 - \rho^2(\hat{\tau})}} \cdot (\tau_q - \hat{\tau})^2 
			\end{aligned}
		\end{equation}
		
		Solving for $(\tau_q - \hat{\tau})^2$
		\begin{equation}
			\begin{aligned}
				(\tau_q - \hat{\tau})^2 &= \frac{R_q(\tau_q) - \frac{2}{\pi} \arcsin(\rho(\hat{\tau}))}{\frac{1}{\pi} \cdot \frac{\rho''(\hat{\tau})}{\sqrt{1 - \rho^2(\hat{\tau})}}} = \frac{\pi \cdot \sqrt{1 - \rho^2(\hat{\tau})} \cdot \left[ R_q(\tau_q) - \frac{2}{\pi} \arcsin(\rho(\hat{\tau})) \right]}{\rho''(\hat{\tau})} 
			\end{aligned}
		\end{equation}
		
		\item Expand $R_q(\tau)$ in a Taylor series around $\tau_0$
		
		\begin{equation}
			\begin{aligned}
				R_q(\tau) &= R_q(\tau_0) + R_q'(\tau_0)(\tau - \tau_0) + \frac{1}{2} R_q''(\tau_0)(\tau - \tau_0)^2  \\
				&= \frac{2}{\pi} \arcsin(\rho(\tau_0)) + \frac{2}{\pi} \cdot \frac{\rho'(\tau_0)}{\sqrt{1 - \rho^2(\tau_0)}} \cdot (\tau - \tau_0)  \\
				&\quad + \frac{1}{2} \cdot \frac{2}{\pi} \left( \frac{\rho''(\tau_0)}{\sqrt{1 - \rho^2(\tau_0)}} + \frac{\rho(\tau_0)[\rho'(\tau_0)]^2}{(1 - \rho^2(\tau_0))^{3/2}} \right) \cdot (\tau - \tau_0)^2 
			\end{aligned}
		\end{equation}
		
		Substituting $\tau_q$ and $\hat{\tau}$ respectively gives
		\begin{equation}
			\begin{aligned}
				R_q(\tau_q) &= R_q(\tau_0) + R_q'(\tau_0)(\tau_q - \tau_0) + \frac{1}{2}R_q''(\tau_0)(\tau_q - \tau_0)^2  \\
				&= \frac{2}{\pi}\arcsin(\rho(\tau_0)) + \frac{2}{\pi}\frac{\rho'(\tau_0)}{\sqrt{1-\rho^2(\tau_0)}} \cdot \Delta\tau_q  \\
				&\quad + \frac{1}{2}\frac{2}{\pi}\left( \frac{\rho''(\tau_0)}{\sqrt{1-\rho^2(\tau_0)}} + \frac{\rho(\tau_0)[\rho'(\tau_0)]^2}{(1-\rho^2(\tau_0))^{3/2}} \right) \cdot \Delta\tau_q^2  \\
			\end{aligned}
		\end{equation}
		
		
		\begin{equation}
			\begin{aligned}
				R_q(\hat{\tau}) &= R_q(\tau_0) + R_q'(\tau_0)(\hat{\tau} - \tau_0) + \frac{1}{2}R_q''(\tau_0)(\hat{\tau} - \tau_0)^2  \\
				&= \frac{2}{\pi}\arcsin(\rho(\tau_0)) + \frac{2}{\pi}\frac{\rho'(\tau_0)}{\sqrt{1-\rho^2(\tau_0)}} \cdot \Delta\tau  \\
				&\quad + \frac{1}{2}\frac{2}{\pi}\left( \frac{\rho''(\tau_0)}{\sqrt{1-\rho^2(\tau_0)}} + \frac{\rho(\tau_0)[\rho'(\tau_0)]^2}{(1-\rho^2(\tau_0))^{3/2}} \right) \cdot \Delta\tau^2 
			\end{aligned}
		\end{equation}
	\end{itemize}
	
	\subsubsection*{D. Questions}
	\begin{enumerate}
		\item According to the Arcsine Law, $R_q(\tau)$ and $\rho(\tau)$ are related by the arcsine function. Since $\arcsin(x)$ is monotonically increasing, when $\rho(\tau)$ reaches its peak, $R_q(\tau)$ should also reach its peak. Does this imply that there is \textit{no bias} between $\tau_q$ and $\hat{\tau}$?
		
		\item The Arcsine Law requires $x(t)$ and $y(t)$ to be continuous-time, scalar, stationary real Gaussian processes. In theory, under this model, $\rho(\tau)$ should be an impulse function (Dirac delta) non-zero only at $\tau = \hat{\tau}$. This would make $\rho(\tau)$ discontinuous and non-differentiable around $\hat{\tau}$, and $\rho''(\hat{\tau})$ would not exist. Does this invalidate the above derivation?
		
		\item Assuming the above issues are resolved, how can we use equations (12), (14), and (15) to further derive an expression for the bias $\delta_\tau$?
	\end{enumerate}
	
	\section*{Method 2: Based on the Definition of Fisher Information}
	
	\subsubsection*{A. Model Assumptions}
	The signal is a known waveform $s(t)$. The received signals are
	\begin{equation}
		\begin{aligned}
			x(t) &= s(t) + n_1(t) \\
			y(t) &= s(t - \tau) + n_2(t) 
		\end{aligned}
	\end{equation}
	
	where $n_1(t), \, n_2(t)$ are zero-mean Gaussian white noise processes. The goal is to estimate the delay $\tau$ over the observation interval $t \in [0, T]$.
	
	\subsubsection*{B. Fisher Information for Unquantized Signal}
	Fisher Information quantifies the identifiability of a parameter estimate and is defined as
	\begin{align}
		\mathcal{I}(\tau) &= \mathbb{E}\left[ \left( \frac{\partial}{\partial \tau} \log p(\mathbf{x}, \mathbf{y}; \tau) \right)^2 \right] 
	\end{align}
	Since the signal $x(t)$ does not contain the parameter $\tau$, we focus on $y(t) = s(t - \tau) + n_2(t)$. Its likelihood function is
	\begin{align}
		p(y(t); \tau) &\propto \exp \left( -\frac{1}{2\sigma^2} \int_{0}^{T} \left[ y(t) - s(t - \tau) \right]^2 dt \right)  \\
		\log p(\mathbf{y}; \tau) &= -\frac{1}{2\sigma^2} \int_{0}^{T} \left[ y(t) - s(t - \tau) \right]^2 dt + C 
	\end{align}
	Differentiating with respect to $\tau$
	\begin{equation}
		\begin{aligned}
			\frac{d}{d\tau} \log p(\mathbf{y}; \tau) &= -\frac{1}{\sigma^2} \frac{d}{d\tau}\left[ \int_{0}^{T} \frac{1}{2} \left[ y(t) - s(t - \tau) \right]^2 dt \right]  \\
			&= \frac{1}{\sigma^2} \int_{0}^{T} \left[ y(t) - s(t - \tau) \right] \cdot s'(t - \tau)  dt 
		\end{aligned}
	\end{equation}
	
	The Fisher Information for the delay $\tau$ is
	\begin{gather}
		\mathcal{I}(\tau) = \mathbb{E}\left[ \left( \frac{d}{d\tau} \log p(\mathbf{y}; \tau) \right)^2 \right] = \frac{1}{\sigma^4} \cdot \mathbb{E}\left[ \left( \int_{0}^{T} \left[ y(t) - s(t - \tau) \right] s'(t - \tau) dt \right)^2 \right]  \\
		y(t) - s(t - \tau) = n_2(t) \Rightarrow \mathbb{E}\left[ \left( \int_{0}^{T} n_2(t) s'(t - \tau) dt \right)^2 \right] = \sigma^2 \int_{0}^{T} [s'(t - \tau)]^2 dt  \\
		\mathcal{I}_{\text{analog}}(\tau) = \frac{1}{\sigma^2} \int_{0}^{T} [s'(t - \tau)]^2 dt 
	\end{gather}
	
	Applying the variable substitution $u = t - \tau$ yields
	\begin{align}
		\mathcal{I}_{\text{analog}} &= \frac{1}{\sigma^2} \int_{-\tau}^{T-\tau} [s'(u)]^2 du \approx \frac{1}{\sigma^2} \int_{-\infty}^{\infty} [s'(u)]^2 du \quad \text{(for $T \gg |\tau|$)} 
	\end{align}
	Here, $s'(u)$ is the first derivative of the signal, representing its slope. This integral can be interpreted as the signal's ``slope energy". A steeper signal slope leads to a lower estimation variance.
	
	\subsubsection*{C. Cross-Correlation under 1-Bit Quantization: Arcsine Law}
	Apply 1-bit quantization to $x(t)$ and $y(t)$
	\begin{align}
		x_q(t) &= \text{sign}(x(t)), \quad y_q(t) = \text{sign}(y(t)) 
	\end{align}
	The cross-correlation of the quantized signals is
	\begin{align}
		R_q(\tau) &= \mathbb{E}[x_q(t)y_q(t-\tau)] 
	\end{align}
	If $x(t)$ and $y(t)$ are zero-mean, jointly Gaussian processes, then according to the Arcsine Law (Van Vleck Formula)
	\begin{align}
		R_q(\tau) &= \frac{2}{\pi} \arcsin(\rho(\tau)) 
	\end{align}
	where $\rho(\tau)$ is the normalized correlation coefficient of the original signals
	\begin{align}
		\rho(\tau) &= \frac{R_{xy}(\tau)}{\sigma_x\sigma_y} = \frac{R_{xy}(\tau)}{\sqrt{R_{xx}(0)R_{yy}(0)}} 
	\end{align}
	
	\subsubsection*{D. Derivative of the Quantized Cross-Correlation}
	To estimate the delay, we are interested in the slope of $R_q(\tau)$
	\begin{align}
		\frac{dR_q(\tau)}{d\tau} &= \frac{2}{\pi} \cdot \frac{1}{\sqrt{1 - \rho^2(\tau)}} \cdot \frac{d\rho(\tau)}{d\tau} 
	\end{align}
	Near $\tau \approx \tau_0$ (assuming $\tau_0=0$ for simplicity), $\rho(\tau) \approx 1$. Defining $\rho(\tau) = 1 - \epsilon(\tau)$, we have $\epsilon(\tau) \approx 0$, and $1 - \rho^2(\tau) \approx 2\epsilon(\tau)$. Thus
	\begin{align}
		\frac{dR_q(\tau)}{d\tau} &\approx \frac{2}{\pi} \cdot \frac{1}{\sqrt{2\epsilon(\tau)}} \cdot \frac{d\rho(\tau)}{d\tau} 
	\end{align}
	
	\subsubsection*{E. Fisher Information Ratio}
	Since the likelihood $\log p(\mathbf{y}_q; \tau)$ for the 1-bit quantized signal cannot be written explicitly, a closed-form expression for the Fisher Information is unavailable. We use an alternative statistic, the cross-correlation function $R_q(\tau)$, and approximate the Fisher Information as the square of the sensitivity of this observable statistic to the parameter, i.e.,
	\begin{align}
		\mathcal{I}_{\text{1-bit}} &\propto \left( \frac{dR_q(\tau)}{d\tau} \right)^2 \approx \left( \frac{2}{\pi} \right)^2 \cdot \frac{1}{2\epsilon(\tau)} \cdot \left( \frac{d\rho(\tau)}{d\tau} \right)^2 
	\end{align}
	Conjecture
	\begin{align}
		\mathcal{I}_{\text{analog}} &\propto \left( \frac{d\rho(\tau)}{d\tau} \right)^2 \\
		\frac{\mathcal{I}_{1-\mathrm{bit}}}{\mathcal{I}_{\text {analog }}}&=\left(\frac{2}{\pi}\right)^{2} \cdot \frac{1}{2 \epsilon(\tau)} 
	\end{align}
	
	\subsubsection*{F. Question}
	The derivation above is not rigorous. However, directly relating $R_q'(\tau)$ to $\rho'(\tau)$ seems insightful for analyzing the performance loss due to 1-bit quantization. Since the Arcsine Law implies $R_q(\tau)$ is simply the arcsine of $\rho(\tau)$, this suggests the \textit{shape} of the cross-correlation function changes (especially when $\rho(\tau) \to 1$), but the \textit{peak location} remains unaffected. Does this imply that 1-bit quantization only increases the Mean Squared Error (MSE) of the estimate without introducing bias?
	
	
\end{document}